% Full local proofs for the qualified Wilks manuscript reconstruction.
\begin{annexitem}{annex:p13.abelard.token}
 {thm:p13-abelard-token}{Realization and content are different premises}
\begin{proof}
Assign \(E_t=0,C_t=1\). Then the stipulated
\(T_t=E_t\land C_t\) is 0, so content alone does not imply truth
of a realized token in this model.
If \(E_t=1\), the conjunction has exactly the value of \(C_t\),
so \(T_t\) and \(C_t\) are equivalent under that premise.
This verifies both branches and identifies the missing assumption.
It does not prove that the conjunction is the complete historical
definition of truth, nor that every candidate token is realized.
\end{proof}
\end{annexitem}

\begin{annexitem}{annex:p13.abelard.square}
 {thm:p13-abelard-square}{Every square relation and the failed O-conversion}
\begin{proof}
Use exactly the extension clauses in the body. Since \(O=\sim A\),
exactly one of A and O is true in two-valued semantics; the same
holds for I and E because \(E=\sim I\).
If A holds, choose \(s\in S\) from its nonemptiness conjunct.
Inclusion gives \(s\in P\), proving I.
If E holds, I fails; A cannot hold since it would imply I.
Therefore O holds. This also proves that A and E cannot both hold.
If I and O both failed, \(\sim O\) would give A, which gives I,
a contradiction. They therefore cannot both fail.

When S is empty, its explicit existence conjunct makes A false,
and \(S\setminus P\) is empty, making the separative O false.
The extinctive O is true as the negation of A.
If S is nonempty, A fails exactly when \(S\not\subseteq P\),
equivalently when some \(s\in S\setminus P\).
Thus O equals the displayed separative O under that condition.
Adding an extra P-nonemptiness presumption does not alter
the two failures when S is empty.

For conversion take distinct \(a,b\), \(S=\{a\}\), \(P=\{a,b\}\).
Both classes are nonempty. \(S\setminus P=\varnothing\),
while \(P\setminus S=\{b\}\).
The two particular negatives have different truth values.
This is a full counterexample even with nonempty predicate
and subject, so existence assumptions alone do not license
the claimed unrestricted O-conversion.
\end{proof}
\end{annexitem}

\begin{annexitem}{annex:p13.abelard.modal}
 {thm:p13-abelard-modal}{Scope, common witnesses, and every restriction direction}
\begin{proof}
Fix a constant domain D and a nonempty set \(S\subseteq D\)
that is the same at every accessible world.
If one accessible world satisfies P for every member of S,
that world witnesses possibility of P for each member.
For failure of the converse, use D and S equal to \(\{a,b\}\),
three worlds \(w,u,v\), and universal accessibility.
Set \(P_w=\varnothing,P_u=\{a\},P_v=\{b\}\).
Each subject possibly satisfies P at w, but no accessible world
satisfies P of both. Thus common possibility is stronger.

Necessity for each member says: for every \(x\in S\)
and every accessible world, \(P(x)\) holds.
Necessary universal predication reverses those two universal
quantifiers. Both statements therefore hold in exactly the
same cases under the fixed-domain/fixed-S hypotheses.
If every member is individually impossible P, no world can
satisfy P of all of them: choose any member of nonempty S
to refute the proposed common world. The converse fails
in the preceding two-subject model, where no common world
exists but both subjects individually have a possible P-world.

For the standing/sitting comparison use one subject a and
worlds w,v with universal accessibility. Let a stand only
at w and sit only at v, and be human at both.
At w a is actually standing and possibly sitting.
At neither world does it both stand and sit, so possible
simultaneous standing/sitting is false.
The boxed implication standing-to-not-sitting is true
at every world, while boxed not-sitting of the actual
standing a is false because a sits at v.
This establishes both opposite scope examples.

For restriction, any witness in nonempty \(W_C\subseteq W\)
is a witness in W. Thus restricted possibility implies
unrestricted possibility. Any property true everywhere
on W is true everywhere on \(W_C\); this proves the
necessity direction.
Let \(W=\{w,v\},W_C=\{w\}\).
Make P false at w and true at v: unrestricted possibility
holds and restricted possibility fails.
Make Q true at w and false at v: restricted necessity
holds and unrestricted necessity fails.
These give both converse failures.
An interval selected by reading illustrates Q; an interval
selected by standing while sitting is excluded illustrates P.
The sets are explicit comparison assumptions, not a proof
that every nature-based or epistemic historical mode is
one of these quantified world operators.
\end{proof}
\end{annexitem}

\begin{annexitem}{annex:p13.abelard.syllogisms}
 {thm:p13-abelard-syllogisms}{Conversions and both indirect reductions in full}
\begin{proof}
For Barbara assume \(M\subseteq P,S\subseteq M\).
Any \(s\in S\) belongs to M and hence P; so \(S\subseteq P\).
If the conclusion is existentially loaded, nonemptiness of S
must also be present, as in the source's affirmative minor.
For Ferio take \(G\cap E=\varnothing\) and witness
\(j\in J\cap G\). It cannot belong to E, so
\(j\in J\setminus E\), proving the particular conclusion.

Negative-universal conversion is symmetry of the empty intersection.
Affirmative-particular conversion is symmetry of a nonempty
intersection, retaining the same witness.
For conversion per accidens, from nonempty \(S\subseteq P\)
choose \(s\in S\). Then \(s\in P\cap S\).
Without the nonempty-S condition, the raw universal inclusion
holds for empty S while that particular conclusion fails.

For Baroco assume \(A\subseteq B\) and choose
\(c\in C\setminus B\). If \(c\in A\), inclusion gives
\(c\in B\), a contradiction. Thus \(c\in C\setminus A\).
For the source's indirect presentation, suppose
\(C\subseteq A\). Barbara gives \(C\subseteq B\),
contradicting the same c-witness. Since C is nonempty
from that witness, failure of the assumed inclusion
indeed supplies a particular C outside A.

For Bocardo assume a witness \(m\in M\setminus P\)
and \(M\subseteq S\). Then \(m\in S\setminus P\)
is the required conclusion.
Indirectly suppose \(S\subseteq P\); Barbara with
\(M\subseteq S\) gives \(M\subseteq P\), contradicting m.
Nonemptiness of M and S is already provided by m.
These are necessary consequences in the declared classical
model. They do not retrospectively change the source's
reported maximum-probability standing or settle every
original separative-negation/existence convention.
\end{proof}
\end{annexitem}

\begin{annexitem}{annex:p13.abelard.middle}
 {thm:p13-abelard-middle}{The exact modal and temporal middle conditions}
\begin{proof}
At a fixed evaluation world, take any x with \(Vx\).
The assertoric minor gives \(Jx\); the major then gives
\(\Diamond Gx\). Universal generalization proves the result.
Replacing the diamond in the major and conclusion by
a box uses precisely the same two actual instantiations.
Likewise, if every E is not-possibly-G and every J is
possibly-G, an x in both would satisfy a proposition
and its negation. Hence J and E are disjoint.

For failure with a modalized minor use domain \(\{a,b\}\),
worlds w,v and universal accessibility.
At w set \(J_w=\{b\},V_w=\{a\},G_w=\varnothing\).
At v set \(J_v=\{a,b\},G_v=\{b\},V_v=\{a\}\).
Every actual J at w is possibly G, witnessed by b at v.
Every actual V at w is possibly J, witnessed by a at v.
Both actual subject classes are nonempty.
But a is G at neither world, so the conclusion that
every actual V is possibly G is false.
The first premise says nothing about a's newly acquired
J-status at v; a global rigidity premise would be extra.

For the tense case use domain \(\{a,b\}\),
CurrentBoy=\(\{b\}\), PastBoy=\(\{a\}\),
PastYoung=\(\{a\}\), CurrentOld=\(\{a\}\).
No current boy belongs to PastYoung, and every current
old person belongs to PastBoy. Yet the current old a
belongs to PastYoung, refuting the proposed conclusion.
PastBoy and CurrentBoy are not the same middle.
An a-history that was a boy, later young and now old,
and a currently young b with no earlier-young stage,
can realize these unary extensions.

For the properly shared harpist case,
Harpist\(\subseteq\)Human\(\subseteq\)FutureDie
gives Harpist\(\subseteq\)FutureDie by the inclusion proof.
For the mortal/living particular case choose a member
of nonempty Mortal, use its inclusions in Living and
FutureDie, and retain it as the witness.
None of these local proofs supplies an unrestricted
historical temporal or modal syllogistic.
\end{proof}
\end{annexitem}

\begin{annexitem}{annex:p13.abelard.warrants}
 {thm:p13-abelard-warrants}{All three transfer cases and the actual-equality countermodel}
\begin{proof}
For the species case, the stated inclusion warrant
\(\forall x(Sx\to Gx)\), instantiated at a, gives
\(Sa\to Ga\). With Sa, detach Ga.
This is a consequence of the supplied warrant, not
proof that arbitrary names are species/genus relatives.

For the tree-definition case assume animalhood implies
the perception condition, and that the tree fails it.
An assumption of its animalhood would give perception,
contradicting that failure; hence it is not animal
under the assumed definition.
For the natural-cause case, instantiate
\(\forall x\forall y((N(x)\land C(x,y))\to N(y))\)
at the named society s and justice j.
The two ground premises \(N(s),C(s,j)\) give the
conjunctive antecedent and then \(N(j)\).
For the proportional-office case instantiate
\(\forall x\forall y((L(x)\land R(x,y))\to L(y))\),
where L is the declared not-selected-by-lot property
and R the asserted office-proportion relation.
With L(pilot) and R(pilot,ruler), detach L(ruler).
In each case every premise used is explicit.
The latter two warrants are not proved by their instances.
For either form, a two-object interpretation with
the property true only of the first object and the
relation connecting it to the second refutes the
conclusion if the transfer premise is omitted.

For actual coextension use worlds w,v, universal
accessibility, and domain \(\{a,b\}\).
At w make Def and Human both \(\{a\}\).
At v make Def=\(\{a,b\}\), Human=\(\{a\}\).
Actual coextension holds, but at v b satisfies Def
without Human. Hence the boxed Def-to-Human
inclusion fails at w.
Human-to-Def holds at both worlds in this model,
showing the two directions can have different
modal standing. This is a comparison countermodel,
not a complete semantics for the original nature
or containment relation.
\end{proof}
\end{annexitem}

\begin{annexitem}{annex:p13.abelard.diagnostic}
 {thm:p13-abelard-diagnostic}{Every operation in the opposite and Alberic chains}
\begin{proof}
Treat \(A\supset B\) as an accepted conditional judgment.
The conclusion is conditional on the operations expressly
listed, not on interpreting this sign as material implication.
For the opposite-topic case assume
\(p\supset\sim q\) and \(q\supset\sim p\).
The source's combination step gives
\((p\land q)\supset(\sim q\land\sim p)\).
The propositional step from the latter conjunction
to \(\sim(p\land q)\), followed by composition,
gives \((p\land q)\supset\sim(p\land q)\).
The propositional step is explicit: if p and q both
held, either supplied negative conjunct would contradict
one of them. Thus it requires an applicable negative
projection/combination, not just a typography change.

For Alberic let \(D=p\land\sim q\) and accept \(p\supset q\).
Projection gives \(D\supset\sim q\).
Contraposition of the accepted conditional gives
\(\sim q\supset\sim p\).
Their composition gives \(D\supset\sim p\).
The other projection gives \(D\supset p\).
Contraposition of that projection gives
\(\sim p\supset\sim D\).
A second composition gives \(D\supset\sim D\).
These are all three initial source judgments and
both subsequent compositions.
Using \(p\land q\supset p\) instead of the required
D-projection would give \(\sim p\supset\sim(p\land q)\),
not the required \(\sim p\supset\sim D\).
This proves why the corrected projection matters.

If the global policy forbids \(X\supset\sim X\)
for every X, each derived result violates it.
At least one accepted premise or one operation's
applicability must therefore be changed.
This proves incompatibility of those specified commitments,
not inconsistency of every historically qualified
connexive theory. The rival school responses locate
different dependencies to restrict.
\end{proof}
\end{annexitem}

\begin{annexitem}{annex:p13.abelard.nesting}
 {thm:p13-abelard-nesting}{The complete truth calculations and their semantic limit}
\begin{proof}
\(K_W(A,B)\) has the same Boolean value at every state:
it records whether all states verify \(A\to B\).
Consequently \(D_W(A,B)\) records coverage of W
by the pointwise truth sets of A and B.
Take W=\(\{h,c,t\}\), with exactly one of Hot,
Cold, Tepid true at its respectively named state.
Then \(D_W(\mathrm{Hot},\mathrm{Cold})=0\),
because t verifies neither; and
\(D_W(\mathrm{Cold},\mathrm{Tepid})=0\),
because h verifies neither.
Therefore the left stronger nesting is
\(D_W(0,\mathrm{Tepid})\), which would require
Tepid everywhere and is false.
The right is \(D_W(\mathrm{Hot},0)\), which would
require Hot everywhere and is also false.

If the inner Hot/Cold combination is the ordinary
pointwise disjunction, its negation is true only at t.
That state satisfies Tepid, so the outer stronger
connection is true. The alternate weak Cold/Tepid
inner combination similarly covers c,t, leaving
only h for the other disjunct. Thus both weakened
repairs verify coverage without identifying their
intensional semantic roles.

For non-associativity on unrestricted inputs use
W=\(\{u,v\}\), A true only at u, and B,C true
only at v. \(D_W(A,B)=1\), so
\(D_W(D_W(A,B),C)=D_W(1,C)=1\).
But \(D_W(B,C)=0\), and
\(D_W(A,D_W(B,C))=D_W(A,0)=0\).
The two associations differ.

For the qualified sickness example take three states:
rock r, well-animal h, not-well-animal s.
Assign Animal values 0,1,1; Well values 0,1,0;
Sick values 0,0,1.
The bare \(K_W(\sim\mathrm{Well},\mathrm{Sick})\)
fails at r. \(K_W(\mathrm{Animal},\mathrm{Sick})\)
fails at h. But the conjunctive restriction
\(K_W(\mathrm{Animal}\land\sim\mathrm{Well},\mathrm{Sick})\)
holds: its antecedent is true only at s, where
the consequent is true.
The two first natural nestings retain one of the
false global inner connections as consequent
with a true outer antecedent, so fail here.

This diagnostic does not reproduce every historical
rejection. In particular its ordinary treatment
of a false antecedent makes
\(K_W(K_W(\mathrm{Animal},\sim\mathrm{Well}),\mathrm{Sick})\)
vacuously true because the inner connection fails at h.
That is precisely why it must not be advertised
as a complete connexive or existence-loaded semantics.
The three assignments are stipulated logical cases,
not clinical facts about all states of health.
\end{proof}
\end{annexitem}

\begin{annexitem}{annex:p13.abelard.explosion}
 {thm:p13-abelard-explosion}{The full five-step conditional argument and its last composition}
\begin{proof}
Let \(X=P\land\sim P\).
Projection gives \(X\supset P\).
Disjunction introduction gives \(P\supset(P\lor Q)\).
Composition gives \(X\supset(P\lor Q)\).
The other projection gives \(X\supset\sim P\).
Combining these two consequences of X gives
the joint premises \(P\lor Q,\sim P\);
disjunctive syllogism gives Q from those premises.
The joint composition therefore gives \(X\supset Q\).
This expands the abbreviated transition from
the source's third and fourth steps to its fifth.
Q is arbitrary, so uniform substitution gives
every claimed arbitrary conclusion.
The interpretation here uses the operations
accepted in the Parvipontanian argument, not
Abelard's stronger conditional plus an unchanged
truth-functional disjunction that he has rejected.
\end{proof}
\end{annexitem}

\begin{annexitem}{annex:p13.abelard.hypothetical-cases}
 {thm:p13-abelard-hypothetical-cases}{All simple conditional/disjunctive variants and the mediate case}
\begin{proof}
Let \(\alpha\) be A or \(\sim A\), and
\(\beta\) be B or \(\sim B\).
For each of the four pairs, the material implication
\(\alpha\to\beta\), with \(\alpha\), gives \(\beta\):
a true antecedent and false consequent would falsify
the supplied implication.
With \(\sim\beta\), that same implication excludes
\(\alpha\), giving \(\sim\alpha\).
If a selected literal is negative, its contradictory
is simplified by classical double-negation elimination.
These four choices and two operations give exactly
the eight displayed conditional cases.

For each of the same four pairs,
\(\alpha\lor\beta,\sim\alpha\) gives \(\beta\),
because a true disjunction with false first disjunct
must have true second disjunct.
Interchanging the disjuncts gives the other case.
Again simplifying double negations gives eight
quality/denial variants.
For the mediate case, from \(\alpha\to\beta\),
\(\beta\to\gamma\), and \(\alpha\), first detach
\(\beta\) and then \(\gamma\).
With \(\sim\gamma\), apply the negative counterpart
first to \(\beta\to\gamma\), obtaining \(\sim\beta\),
and then to \(\alpha\to\beta\), obtaining \(\sim\alpha\).
Every premise for every step is supplied.
These proofs cover the stated local classical
comparison, not the full 128-case original taxonomy
or its differently interpreted temporal compounds.
\end{proof}
\end{annexitem}
